\lysh{38856}{4}

\begin{alist}
\item
\begin{rlist}
\item Αν αγοράσει ποδήλατο και το χρησιμοποιήσει για $5$ μήνες, θα πληρώσει:
$400 + 5 \cdot 10 = 450\text{\euro}$.
\item Αν νοικιάσει ποδήλατο και το χρησιμοποιήσει για $5$ μήνες, θα πληρώσει:
$35 \cdot 5 = 175\text{\euro}$.
\end{rlist}

\item Το κόστος της αγοράς του καινούριου ποδηλάτου μαζί με τα μηνιαία έξοδα συντήρησης για $x$ μήνες δίνονται από τη σχέση: $f(x) = 400 + 10x$, όπου $x \in \mathbb{N}$, ενώ το κόστος ενοικίασης ποδηλάτου για $x$ μήνες είναι: $g(x) = 35x$, $x \in \mathbb{N}$.

\item Η αγορά ενός καινούριου ποδηλάτου μαζί με τη μηνιαία συντήρησή του συμφέρουν πιο πολύ τον Στέφανο όταν:
\[400 + 10x < 35x \Leftrightarrow 400 < 35x - 10x \Leftrightarrow 400 < 25x \Leftrightarrow x > \dfrac{400}{25} \Leftrightarrow x > 16.\]
Επομένως, η αγορά ποδηλάτου συμφέρει περισσότερο τον Στέφανο από τη μηνιαία ενοικίαση, εάν έχει σκοπό να μετακινείται με ποδήλατο για τουλάχιστον $17$ μήνες.

\textit{Εναλλακτικός τρόπος:}
$400 + 10x = 35x \Leftrightarrow 400 = 35x - 10x \Leftrightarrow 400 = 25x \Leftrightarrow x = \dfrac{400}{25} \Leftrightarrow x = 16$. Για $x = 16$ και οι δύο επιλογές συμφέρουν το ίδιο τον Στέφανο, ενώ αν χρησιμοποιήσει το ποδήλατο για περισσότερους από $16$ μήνες (δηλαδή $17$) τον συμφέρει η αγορά.

\item
\begin{rlist}
\item Αν επιλέξει να αγοράσει ποδήλατο:
\[400 + 10x \le 600 \Leftrightarrow 10x \le 600 - 400 \Leftrightarrow 10x \le 200 \Leftrightarrow x \le \dfrac{200}{10} \Leftrightarrow x \le 20.\]
Για να μην ξεφύγει από τον οικονομικό του προϋπολογισμό, θα πρέπει να χρησιμοποιήσει το ποδήλατο που θα αγοράσει το πολύ μέχρι $20$ μήνες.

Αν επιλέξει να νοικιάσει ποδήλατο:
\[35x \le 600 \Leftrightarrow x \le \dfrac{600}{35} \Leftrightarrow x \le 17,14.\]
Εφόσον $x \le 17,14$, αν νοικιάσει το ποδήλατο για $18$ μήνες ή περισσότερους, θα ξεφύγει από τον προϋπολογισμό του. Άρα για να μην ξεφύγει από τον οικονομικό του προϋπολογισμό, θα πρέπει να το νοικιάσει μέχρι και $17$ μήνες.
\item Σύμφωνα με το προηγούμενο ερώτημα, είτε με αγορά είτε με ενοικίαση, τα $600\text{\euro}$ καλύπτονται σε περισσότερους από $12$ μήνες. Άρα και με τις δύο επιλογές θα είναι εντός του ετήσιου προϋπολογισμού του.
\end{rlist}
\end{alist}
